\chapter{Napredni koncepti objektno orijentiranog programiranja}


\section{Statički članovi klasa}

Uobičajeni članovi klase pripadaju svakom objektu zasebno. 
Primjerice, svaki objekt sljedeće klase
\begin{minted}{C++}
class Duznik {
    double iznos, tecaj;
};
\end{minted}
ima vlastitu kopiju člana \texttt{iznos} i \texttt{tecaj}.
Međutim, ponekad želimo da podatak bude zajednički svim objektima određene klase. 
U takvim situacijama koristimo \textbf{statičke članove klase} (engl.\ \textit{static members}).
Statički član deklarira se pomoću ključne riječi 
\texttt{static}. 
Za razliku od običnih članova, postoji samo jedna njegova kopija u memoriji, neovisno o broju stvorenih objekata.
Ako izmijenimo klasu \texttt{Duznik} iz gornjeg primjera
na sljedeći način:
\begin{minted}{C++}
class Duznik {
    double iznos;
    static double tecaj;
};
\end{minted}
tada za sve objekte tipa \texttt{Duznik}
postoji samo jedan \texttt{tecaj}
(neovisno o broju stvorenih objekata tipa \texttt{Duznik},
uključujući i slučaj kada nije stvoren niti jedan objekt tog tipa).
Iako nisu dijelovi pojedinih objekata, možemo preko njih
pristupati statičkim članovima pripadne klase.
Primjerice, ako izmijenimo kod klase \texttt{Duznik} tako da \texttt{tecaj}
bude deklariran kao javan
(a ne privatan) dio klase,
obje sljedeće upotrebe tečaja su ispravne:
\begin{minted}{C++}
Duznik::tecaj = 0.95;
Duznik duznik;
duznik.tecaj = 0.99;
\end{minted}
\medskip

Budući da statički član ne pripada pojedinom objektu nego cijeloj klasi, potrebno ga je definirati izvan klase.
Primjerice, za klasu \texttt{Duznik} iz primjera to bismo 
mogli napraviti ovako:
\begin{minted}{C++}
double Duznik::tecaj = 1.56;
\end{minted}

Klasa može sadržavati i \textbf{statičke funkcije}. 
Takve funkcije deklariraju se s ključnom riječi \texttt{static}
te također pripadaju klasi, a ne pojedinom objektu.
Zbog toga statička funkcija može pristupati samo statičkim članovima klase.


\begin{zadatak}\label{zad:racuni-static}
    Izmijenite klasu \texttt{Racun} iz zadatka \ref{zad:racuni}
    tako da član klase koji se odnosi na tečaj ne pripada pojedinom
    računu nego postoji samo jedan za sve račune.
    Kako bi se podatak o tečaju (kojem korisnik nema pristup jer 
    je on privatan) mogao mijenjati izvan klase,
    dodajte u sučelje klase metodu
    \texttt{void promijeni(double novi\_tecaj)}.
    Ta metoda postavlja vrijednost tečaja na
    vrijednost \texttt{novi\_tecaj}
    samo ako je ta vrijednost u intervalu 
    $[0.9,1]$.

    Napišite i primjer programa za testiranje koji
    treba postaviti vrijednost tečaja 
    i prije deklaracije bilo kojeg objekta tipa \texttt{Racun}.
\end{zadatak}


\begin{rjesenje}
    Postavimo \texttt{tecaj} kao statički član klase \texttt{Racun}
    te deklariramo javnu statičku metodu \texttt{promijeni}:
\begin{minted}[
        linenos,
        numbersep=5pt
    ]{C++}
class Racun {
private:
    std::string ime;
    double iznos;
    static double tecaj;
    double pretvori();
public:
    Racun(double, std::string im = "nepoznat");
    void info();
    static void promijeni(double);
};
\end{minted}
Pritom smo izmijenili i konstruktor klase koji sada više ne treba primiti vrijednost tečaja obzirom da on ne pripada
    pojedinom objektu koji se konstruira.
Izmijenimo prema tome implementaciju konstruktora klase u \textit{Racun.cpp} datoteci na sljedeći način:
\begin{minted}[
        linenos,
        numbersep=5pt
    ]{C++}
Racun::Racun(double iz, std::string im)
    : ime(im), iznos(iz) {}
\end{minted}
Kako tečaj sad nije dio pojedinog računa pa
nije definiran kad konstruiramo neki račun,
definiramo ga i inicijaliziramo izvan klase,
u \textit{Racun.cpp} datoteci:
\begin{minted}[
        linenos,
        numbersep=5pt
    ]{C++}
double Racun::tecaj = 0.94;
\end{minted}
Ovdje smo vrijednost $0.94$ odabrali na slučajan način.
\medskip

Preostaje još implementirati statičku metodu \texttt{promijeni} u \textit{Racun.cpp} datoteci:
\begin{minted}[
        linenos,
        numbersep=5pt
    ]{C++}
void Racun::promijeni(double novi) {
    if(novi >= 0.9 && novi <= 1)
        tecaj = novi;
}
\end{minted}
Uočite da u gornjem kodu nismo ponovo navodili ključnu riječ \texttt{static}.
\medskip

Primjer traženog programa za testiranje:
\begin{minted}[
        linenos,
        numbersep=5pt
    ]{C++}
Racun::promijeni(0.99);
Racun r1(100.99,"John"), r2(123.23);
r1.promijeni(0.95); // ili Racun::promijeni(0.95);
r1.info();
r2.info();
\end{minted}
\end{rjesenje}

\begin{razmisljanje}
    \begin{enumerate}
        \item 
        Objasnite na primjeru programa za testiranje iz
        rješenja prethodnog zadatka,
        do kojeg problema bi došlo kada bismo metodu \texttt{promijeni}
        deklarirali kao nestatičku metodu klase \texttt{Racun}.
        
        \item 
        Ako je poznato da se objekti mogu definirati samo jednom, 
        obrazložite zašto definiciju i inicijalizaciju
        tečaja:
        \begin{minted}[
        linenos,
        numbersep=5pt
    ]{C++}
double Racun::tecaj = 0.94;
\end{minted}
    nismo mogli umjesto u \texttt{Racun.cpp} datoteku
    staviti u \texttt{Racun.h} datoteku.

    \item 
    Objasnite zašto sljedeća verzija implementacije metode \texttt{promijeni}
    klase \texttt{Racun} iz prethodnog zadatka nije ispravna
    (u smislu da je nije moguće kompajlirati):

    \begin{minted}[
        linenos,
        numbersep=5pt
    ]{C++}
void Racun::promijeni(double novi) {
    cout << iznos;
    if(novi >= 0.9 && novi <= 1)
        tecaj = novi;
}
    \end{minted}
\end{enumerate}
\end{razmisljanje}



Upotrebom statičkih članova klase možemo u svakom trenutku
imati evidenciju ukupnog broja objekata tipa te klase
koji postoje u programu.


\begin{zadatak}\label{zad:racuni-ukupno}
    Izmijeniti klasu \texttt{Racun} iz zadatka \ref{zad:racuni-static}
    tako da klasa ima član koji sadrži podatak
    o trenutnom ukupnom broju objekata te klase u programu.
    Kako bi korisnik mogao dobiti podatak o tome broju,
    dodati klasi metodu naziva \texttt{koliko}
    koja vraća korisniku traženi broj.
    Napišite i program za testiranje koji sadrži 
    (globalnu) funkciju
    \texttt{void funkcija(Racun r)}
    koja ispisuje podatke o računu \texttt{r}
    i ukupan broj računa u trenutku izvršavanja te funkcije.
\end{zadatak}


\begin{rjesenje}
    Dodamo klasi \texttt{Racun} privatan statički brojač.
    Važno je da je taj brojač privitan kako korisnik
    ne bi (ne)namjerno mogao izravno promijeniti njegovu vrijednosti.
    Također dodajemo statičku metodu \texttt{static int koliko()}
    za dohvat vrijednosti brojača u sučelju.
    Vrijednost brojača povećavat ćemo u konstruktoru i konstruktoru kopiranjem obzirom da 
    se pomoću njih stvaraju novi računi
    (i novi računi koji su kopije postojećih računa),
    a smanjivati u destruktoru (jer tada nestaju računi).
    Prema tome, dodajemo još i konstruktor kopiranjem i destruktor.

    \begin{minted}[
        linenos,
        numbersep=5pt
    ]{C++}
class Racun {
private:
    static int broj_racuna;
    ...
public:
    static int koliko();
    Racun(const Racun&);
    ~Racun();
    ...
};
    \end{minted}
U datoteku \textit{Racun.cpp} dodajemo definiciju brojača
(obzirom da je on statički član klase)
te ga početno inicijaliziramo na $0$.
U konstruktoru dodajemo povećavanje brojača,
te dodajemo konstruktor kopiranjem i destruktor koji povećavaju ili smanjuju brojač.
Naposlijetku dodajemo i implementaciju metode za dohvat vrijednosti brojača.

\begin{minted}[
        linenos,
        numbersep=5pt
    ]{C++}
int Racun::broj_racuna = 0;

Racun::Racun(double iz, std::string im)
    : ime(im), iznos(iz) {
        ++broj_racuna;
}

Racun::Racun(const Racun& r)
    : ime(r.ime), iznos(r.iznos) {
        ++broj_racuna;
}

Racun::~Racun() {
    --broj_racuna;
}

int Racun::koliko() {
    return broj_racuna;
}
\end{minted}

Sadržaja datoteke \texttt{main.cpp} s 
primjerom traženog programa za testiranje:

\begin{minted}[
        linenos,
        numbersep=5pt
    ]{C++}
#include <iostream>
#include "Racun.h"

using namespace std;

void funkcija(Racun r) {
    r.info();
    cout << r.koliko() << endl;
}

int main(void) {
    cout << Racun::koliko() << endl;
    Racun::promijeni(0.99);
    Racun r1(100.99,"John"), r2(123.23);
    r1.promijeni(0.95);
    funkcija(r1);
    funkcija(r2);
    cout << Racun::koliko() << endl;
    return 0;
}
\end{minted}
\end{rjesenje}

\begin{razmisljanje}
    \begin{enumerate}
        \item
        Što biste izmijenili u rješenju prethodnog zadatka
        kako bi imali evidenciju i o ukupnom broju
        objekata tipa \texttt{Racun}
        koji su stvoreni kao kopije nekog drugog objekta
        tog tipa?
        Kako biste u destruktoru znali kad se uništava 
        takva kopija od uništavanja računa koji nije nastao kao kopija?
    \end{enumerate}
\end{razmisljanje}


\section{\texttt{this} pokazivač}

Funkcije članice klase pristupaju objektu tipa te klase
na kojem su pozvane kroz dodatan implicitan parametar koji se naziva \texttt{\textbf{this}}.
Pri pozivu funkcije, pokazivač \texttt{this} inicijalizira se
adresom objekta na kojem je funkcija pozvana.
\medskip

Primjerice, promotrimo metodu \texttt{info} klase \texttt{Racun} iz zadatka \ref{zad:racuni}:
\begin{minted}{C++}
void Racun::info() {
    cout << ime << endl
         << '\t' << iznos << "EUR" << endl
         << '\t' << pretvori() << "CHF" << endl;
}
\end{minted}
Pri pozivu te metode u sljedećem kodu
\begin{minted}{C++}
Racun jedan_racun(100);
jedan_racun.info();
\end{minted}
\texttt{this} pokazivač koji možemo koristiti unutar
metode \texttt{Racun::info} sadrži adresu objekta \texttt{jedan\_racun},
dok dereferenciranjem tog pokazivača,
tj.\ izrazom \texttt{*this},
dobivamo upravo objekt \texttt{jedan\_racun}.

Prema tome, našu metodu \texttt{info} mogli smo napisati i ovako
(dobiva se potpuno ista metoda):
\begin{minted}{C++}
void Racun::info() {
    cout << this->ime << endl
         << '\t' << this->iznos << "EUR" << endl
         << '\t' << this->pretvori() << "CHF" << endl;
}
\end{minted}

\texttt{this} parametar metode za nas je definiran implicitno.
Drugim riječima,
poziv metode izrazom 
\texttt{jedan\_racun.info();} i gornju metodu \texttt{info} kompajler kao da prevede u
\begin{center}
    \texttt{Racun::info(\&jedan\_racun);}
\end{center}
i
\begin{center}
    \texttt{void Racun::info(Racun* const this) \{...\}}
\end{center}




\section{\texttt{const} objekti, pokazivači, reference
i funkcije članice}


Ako želimo definirati varijablu čija se vrijednost nakon inicijalizacije ne može mijenjati (ni namjerno ni slučajno tijekom izvođenja programa), prilikom deklaracije koristimo ključnu riječ \texttt{const}.
Obzirom da se vrijednost takve varijable nakon inicijalizacije ne može
mijenjati, ona mora biti inicijalizirana.
Primjerice, ako deklariramo
\begin{minted}{C++}
const double t = 7.5345;
\end{minted}
tada možemo, primjerice, ispisati vrijednost te varijable
pomoću izraza \texttt{cout <\!< t}
ali joj ne možemo tu vrijednost promijeniti pomoću izraza \texttt{t *= 2}.
\medskip

Ključna riječ \texttt{const} može se koristiti i prilikom stvaranja objekata klase. 
Nakon inicijalizacije takvih objekata nije dopušteno mijenjati vrijednosti njihovih članova.
Primjerice, možemo napisati
\begin{minted}{C++}
const Racun primjer(100,"test");
\end{minted}
No, ne možemo funkciji koja prima ne-\texttt{const} pokazivač ili 
ne-\texttt{const} referencu
na račun, poput funkcije s prototipom
\begin{minted}{C++}
bool jel_bogat(Racun& racun);
\end{minted}
proslijediti adresu računa ili račun \texttt{primjer}
jer bi ta funkcija mogla preko pokazivača
ili reference pokušati promijeniti taj račun
(što predstavlja grešku obzirom da je \texttt{primjer} označen sa \texttt{const}).
Stoga je potrebno koristiti pokazivač koji pokazuje na \texttt{const}
ili \texttt{const} referencu.
U našem primjeru, funkciju \texttt{jel\_bogat}
izmijenili bismo tako da ima sljedeći prototip:
\begin{minted}{C++}
bool jel_bogat(const Racun& racun);
\end{minted}
Drugim riječima, funkcija time jamči da neće mijenjati objekt na koji 
pokazuje ili koji referencira \texttt{racun}.
Također, u toj funkciji na takvim objektom ne mogu se
pozivati druge funkcije ili metode koje ne jamče da se takav objekt neće pokušati promijeniti.
Istaknimo da se takvim funkcijama i metodama mogu prosljeđivati
objekti ili pokazivači na objekte koji nisu označeni sa \texttt{const}
jer upotreba \texttt{const} u funkcijama i metodama koju
smo opisali samo jamči da ih funkcija ili metoda neće mijenjati
(oni se i dalje mogu mijenjati na druge načine).
\medskip

Metode koje pozivamo nad objektom označenim sa \texttt{const},
također moraju jamčiti da neće pokušati promijeniti taj objekt.
Primjerice,
kako bi za račun \texttt{primjer} iz gornjeg primjera napravili poziv neke metode \texttt{zbroji}, poput
\begin{minted}{C++}
double zbroj = primjer.zbroji();
\end{minted}
potrebno je iza liste parametara metode dodati
(i u deklaraciji i pri implementaciji metode) ključnu riječ \texttt{const}:
\begin{minted}{C++}
double primjer() const;
\end{minted}
Dodavanjem \texttt{const} na kraj metode implicitni pokazivač \texttt{this} postaje tipa 
\begin{center}
    \texttt{const Racun* const this}
\end{center}
(čitajući zdesna nalijevo: \texttt{this} je konstantni
pokazivač na konstantni objekt klase \texttt{Racun}, tj.\ ne može se promijeniti da pokazuje na neki drugi objekt
te objekt na koji pokazuje nije moguće promijeniti). 
Time metoda jamči da se objekt koji tu metodu koristi neće pokušati mijenjati.


\begin{savjet}
U pravilu je dobro stavljati \texttt{const} za svaki argument
funkcije i metode za koji znamo da ga funkcija ili metoda ne treba mijenjati. 
To se naziva \textit{const-correctness} i smatra se dobrom praksom u programiranju.
Tako da ćemo u nastavku pokušati pri pisanju funkcija i metoda što više koristiti \texttt{const}
(na mjestima gdje je to moguće).
\end{savjet}


\begin{zadatak}
    Izmijeniti klasu \texttt{Racun} te (globalnu) funkciju \texttt{funkcija}
    iz rješenja zadatka \ref{zad:racuni-ukupno}:
    \begin{minted}[
        linenos,
        numbersep=5pt
    ]{C++}
void funkcija(Racun r) {
    r.info();
    cout << r.koliko() << endl;
}
    \end{minted}
    tako da ta funkcija upotrebom pokazivača ne dovodi niti do jednog kopiranja
    nekog objekta tipa \texttt{Racun}
    te da se sljedeći kod (primjerice, u \texttt{main} funkciji) uspješno kompajlira:
    \begin{minted}[
        linenos,
        numbersep=5pt
    ]{C++}
const Racun primjer(100,"test");
funkcija(&primjer);
    \end{minted}
\end{zadatak}


\begin{rjesenje}
    Izmijenimo funkciju \texttt{funkcija} tako da joj je argument pokazivač i 
    to pokazivač na \texttt{const} objekt
    (kako bi mogli imati poziv funkcije i za \texttt{const} objekte kao što je
    objekt \texttt{primjer} iz teksta zadatka):
\begin{minted}[
        linenos,
        numbersep=5pt
    ]{C++}
void funkcija(const Racun * r) {
    r -> info();
    cout << r -> koliko() << endl;
}
\end{minted}

Budući da je parametar \texttt{r} tipa \texttt{const Racun*}, 
preko njega mogu se pozivati samo metode označene s \texttt{const}. 
Prema tome potrebno je dodati \texttt{const}
iza liste parametara metode \texttt{info}
kako bi naznačili da ona neće pokušati promijeniti objekt nad kojim se poziva (u ovom slučaju, 
objekt na koji pokazuje pokazivač \texttt{r}).
Dodamo \texttt{const} u \textit{Racun.h} datoteci pri deklaraciji \texttt{info} metode:
\begin{minted}[
        linenos,
        numbersep=5pt
    ]{C++}
void info() const ;
\end{minted}
te u \textit{Racun.cpp} datoteci pri implementaciji te metode:
\begin{minted}[
        linenos,
        numbersep=5pt
    ]{C++}
void Racun::info() const {
    cout << ime << endl
         << '\t' << iznos << "EUR" << endl
         << '\t' << pretvori() << "CHF" << endl;
}
\end{minted}
Budući da je sada \texttt{info} konstantna metoda, 
unutar nje se mogu pozivati samo konstante metode. 
Zato i \texttt{pretvori} mora biti deklarirana s \texttt{const} 
(što je u redu jer niti ne želimo da ona mijenja objekt nad kojim se poziva).
Dodamo \texttt{const} u \textit{Racun.h} datoteci pri deklaraciji \texttt{pretvori} metode:
\begin{minted}[
        linenos,
        numbersep=5pt
    ]{C++}
double pretvori() const ;
\end{minted}
te u \textit{Racun.cpp} datoteci pri implementaciji te metode:
\begin{minted}[
        linenos,
        numbersep=5pt
    ]{C++}
double Racun::pretvori() const {
    return iznos * tecaj;
}
\end{minted}
Istaknimo da funkcija \texttt{funkcija}
koristi poziv statičke metode \texttt{koliko}
koju ne moramo mijenjati.
To je zato što statičke funkcije članice ne pripadaju pojedinom
objektu tipa klase pa nemaju implicitno \texttt{this} pokazivač. 
Prema tome se niti ne mogu deklarirati kao \texttt{const} (niti koristiti \texttt{this} pokazivač u svojoj
implementaciji).
\end{rjesenje}

\begin{razmisljanje}
    \begin{enumerate}
        \item 
        Zašto se sljedeći kod
        \begin{minted}{C++}
const Racun primjer(100,"test");
funkcija(primjer);
        \end{minted}
        uspješno kompajlira ako ne promijenimo funkciju \texttt{funkcija}
        nego ona ostane ista kao u rješenju zadatka \ref{zad:racuni-ukupno}? 
        \item 
        Obrazložite zašto se sljedeći kod uspješno kompajlira:
        \begin{minted}{C++}
void funkcija(const Racun *r) {
    Racun novi(1.23,"abc");
    r = &novi;
    r -> info();
    cout << r -> koliko() << endl;
}
    \end{minted}

    \item 
    Obrazložite zbog koje se dvije linije sljedeći kod ne može uspješno kompajlirati:
    \begin{minted}[
        linenos,
        numbersep=5pt
    ]{C++}
void Racun::info() const {
    const Racun* const b;
    Racun n(1,"asd");
    this = &n;
    cout << ime << endl;;
}
\end{minted}
    \end{enumerate}
\end{razmisljanje}


\section{Upotreba \texttt{this} pokazivača}


\subsection{Ulančavanje poziva metoda}

Pokazivač \texttt{this} može se koristiti za ulančavanje poziva metoda (engl. \emph{method chaining}). Takav način korištenja već smo susreli kod klase \texttt{std::string}. Na primjer, za objekt \texttt{rijec} tipa te klase moguće je napisati:

\begin{minted}{C++}
rijec.append("je").append("dan");
\end{minted}

Metoda \texttt{append} deklarirana je na sljedeći način:

\begin{minted}{C++}
string& std::string::append(const string& str);
\end{minted}

Budući da metoda vraća referencu na objekt nad kojim je pozvana, moguće je odmah pozvati novu metodu nad rezultatom prethodnog poziva. U implementaciji se to najčešće ostvaruje vraćanjem izraza \texttt{*this}. Stoga je drugi poziv metode \texttt{append} u prethodnom primjeru ekvivalentan pozivu

\begin{minted}{C++}
rijec.append("dan");
\end{minted}

nad istim objektom \texttt{rijec}. Na taj način omogućeno je ulančavanje više uzastopnih poziva metoda u jednom izrazu.


\begin{zadatak}
    Dodajte klasi \texttt{Racun} 
    iz zadatka \ref{zad:racuni-ukupno} 
    funkciju \texttt{uplati} za uplatu na pojedini račun. 
    Funkcija prima iznos (tipa \texttt{double}) za uplatu.
    \begin{itemize}
        \item[(a)]
        Treba li ta funkcija biti statička ili ne-statička?
        \item[(b)]
        Omogućite ulančavanje poziva te funkcije, poput
        \begin{center}
            \texttt{r1.uplati(12).uplati(0.23);}
        \end{center}
        pri čemu je \texttt{r1} tipa \texttt{Racun}.
    \end{itemize}
\end{zadatak}


\begin{rjesenje}
    \begin{itemize}
        \item[(a)]
        Funkcija \texttt{uplati} treba biti nestatička, jer mijenja stanje konkretnog objekta klase \texttt{Racun}.

        \item[(b)]
        Dodamo u sučelje klase \texttt{Racun}
        potrebnu deklaraciju (u \texttt{Racun.h} datoteci):
        \medskip

        \begin{minted}{C++}
class Racun {
    ...
public:
    Racun& uplati(double);
    ...
};
        \end{minted}
        \medskip

        Dodamo u \texttt{Racun.cpp} datoteku:
        \medskip

        \begin{minted}{C++}
Racun& Racun::uplati(double i) {
    iznos += i;
    return *this;
}
        \end{minted}
    \end{itemize}
\end{rjesenje}


\begin{razmisljanje}
\begin{enumerate}
    \item
    U prethodnom rješenju, koju bi grešku pri kompajliranju dobili ako bi funkcija \texttt{uplati} bila statička?
    \item
    Zašto je povratni tip \texttt{Racun\&}, a ne \texttt{Racun}?
    Koji bi tada bio rezultat izvršavanja sljedećeg koda
    (za objekt \texttt{r1} tipa \texttt{Racun}):
    \begin{minted}{C++}
r1.uplati(100).uplati(200);
    \end{minted}
    \item 
    Zašto funkcija \texttt{uplati} vraća \texttt{*this}, a ne samo \texttt{this}?
    Kako bi izgledala funkcija (te upotreba te funkcije) 
    koja bi umjesto \texttt{*this} vraćala \texttt{this}?
    \item 
    Koji od sljedeća dva retka je ispravan obzirom na trenutni kod koji imamo (ako je \texttt{r1} objekt klase \texttt{Racun}):
    \begin{minted}{C++}
r1.uplati(100).info();
r1.info().uplati(100);
    \end{minted}
    Što bi trebalo promijeniti kako bi se oba gornja retka kompajlirala?
\end{enumerate}
\end{razmisljanje}


\subsection{Vođenje evidencije postojećih objekata}

Želimo napisati funkciju koja omogućuje dohvat informacija o svim objektima klase \texttt{Racun} koji u nekom trenutku postoje u programu. 
Budući da takva funkcionalnost nije vezana uz pojedini objekt klase \texttt{Racun}, već uz klasu u cjelini, implementirat ćemo je kao statičku funkciju.
\medskip

Kako bismo mogli pristupiti svim postojećim računima, klasa će sadržavati statičku listu pokazivača na objekte klase \texttt{Racun}. 
Prilikom stvaranja novog objekta, pokazivač \texttt{this} dodaje se u tu listu, čime se svaki račun automatski registrira među aktivnim računima.
Slično, uklanjamo pokazivač \texttt{this} pri uništavanju
pojedinog objekta.
\medskip

Na taj način statičke funkcije klase mogu pristupiti svim postojećim objektima i, primjerice, ispisati njihove podatke ili izračunati ukupno stanje svih računa.

\begin{zadatak}
    Dodajte klase \texttt{Racun} listu odgovarajućeg tipa tako da
    klasa \texttt{Racun} omogućuje pristup svim objektima tipa te klase koji trenutno postoje u programu.
    Dodajte klase i (statičku) funkciju \texttt{void info\_svi()}
    koja će ispisati informacije o svim računima
    (pozivom \texttt{info} metode pojedinog računa).
\end{zadatak}


\begin{rjesenje}
    Dopunimo deklaraciju klase \texttt{Racun} (u \texttt{Racun.h} datoteci):

    \begin{minted}[
        linenos,
        numbersep=5pt
    ]{C++}
#include <list>
class Racun {
private:
    ...
    static std::list<Racun*> svi;
public:
    static void info_svi();
    ...
};
    \end{minted}

    Uočite da smo dodali \texttt{\#include <list>} u datoteku \texttt{Racun.h}, a također ćemo to sad dodati i u 
    datoteku \texttt{Racun.cpp}.
    U toj \texttt{Racun.cpp} datoteci dopunit ćemo konstruktor, konstruktor kopiranja te destruktor
    kako bi lista svih objekata tipa \texttt{Racun}
    u svakom trenutku bila ažurna.

    \begin{minted}[
        linenos,
        numbersep=5pt
    ]{C++}
#include <list>
Racun::Racun(double iz, std::string im) : ... {
    ...
    svi.push_back(this);
}
Racun::~Racun() {
    ...
    svi.remove(this);
}
Racun::Racun(const Racun& r) : ... {
    ...
    svi.push_back(this);
}
    \end{minted}

    Dodamo definiciju statičke liste u datoteku \texttt{Racun.cpp} zajedno s implementacijom
    funkcije \texttt{Racun::info\_svi} koja pomoću
    iteratora prolazi listom pokazivača na sve račune i 
    ispisuje informacije o pojedinom računu
    (pozivom pripadne funkcije \texttt{info} pojedinom računa).

    \begin{minted}[
        linenos,
        numbersep=5pt
    ]{C++}
list<Racun*> Racun::svi;
void Racun::info_svi() {
    cout << "Informacije o svim racunima:" << endl;
    list<Racun*>::iterator it;
    for(it = svi.begin(); it != svi.end(); ++it)
        (*it)->info();
}
    \end{minted}
\end{rjesenje}


\begin{razmisljanje}
\begin{enumerate}
    \item
    Koja je u prethodnom rješenju prednost upotrebe liste u odnosu na, primjerice, vektor?
    \item
    Zašto umjesto liste pokazivača na račune ne možemo koristiti listu računa?
    \item 
    Ako su \texttt{primjer}, \texttt{r1} i \texttt{r2}
    jedini računi koji se pojavljuju u programu,
    što će ispisati funkcija \texttt{info\_svi} u svakom od sljedećih poziva:

    \begin{minted}{C++}
Racun::info_svi();
const Racun primjer(100,"test");
{
    Racun r1(100.99,"John"), r2(123.23,"Mary");
    r1.uplati(100).uplati(321.23);
    Racun::info_svi();
}
Racun::info_svi();
    \end{minted}
\end{enumerate}
\end{razmisljanje}